\scanpage{31}
\handwritten{11.9.1967}\\

Dear Anantharaman,\\

Matsumura proved that if $X$ is proper over a field $k$, then $\underline{\operatorname{Aut}}_{X/k}$ is representable by a group scheme locally of finite type over $k$. I think I can systematize the key step of his argument in the following way. Consider a scheme $S$, and a morphism

\[
\phi: Z\handwritten{\longrightarrow} X
\]

of $S$-schémes which are proper, flat and of finite presentation. Let $Y$ be locally of finite présentation and separated over $S$, then $\phi$ induces a homomorphisme of functors $u\handwritten{\leadsto}u\phi$ ;

\[
\phi':\underline{\operatorname{Hom}}_S(X,Y)\handwritten{\longrightarrow}\underline{\operatorname{Hom}}_S(Z,Y)\quad .
\]

Then one can défine a subfunctor of $\underline{\operatorname{Hom}}_S(\deleted{\unsure{Z}}X,Y)$ where $\phi'$ is ``unramified'' in a rather obvious sense, and this turns out to be an ``open subfunctor'', say $\underline{\operatorname{Hom}}_S(X,Y;\phi)$ . Now look at the induced homomorphism

\[
\underline{\operatorname{Hom}}_S(X,Y;\phi)\handwritten{\longrightarrow}\underline{\operatorname{Hom}}_S(Z,Y)\,.
\]

Using the main result of Murre's talk, one can prove that the latter morphism is representable by unramified separated morphisms locally of finite presentation; as a consequence, if $\underline{\operatorname{Hom}}_S(Z,Y)$ is representable, so is $\underline{\operatorname{Hom}}_S(X,Y;\phi)$.

To get, given $X$ and $Y$, a representability theorem for $\underline{\operatorname{Hom}}_S(X,Y)$, on\deleted{\unsure{ye|es|ge}}\unsure{e} tries to find morphisms $\phi_i:Z_i\handwritten{\longrightarrow}X$ as above, such that the open subfunc\deleted{\unsure{ot}}tors $\underline{\operatorname{Hom}}_S(X,Y;\phi_i)$ cover $\underline{\operatorname{Hom}}_S(X,Y)$ (as a fpqc sheaf), and such that the functor\deleted{\unsure{z}}s $\underline{\operatorname{Hom}}_S(\deleted{\unsure{X}}Z_i,Y)$ are all representable. If for instance $\deleted{\unsure{X}}S$ is the spectrum of a field $k$, and if $X$ has ``enough'' points radicial over $k$ (which is always true if $k$ is alg. closed) then we can take for $Z_i$ all finite subs\deleted{\unsure{h}}\unsure{c}hemes of $X$ whose points are radicial over $k$, and we get

% archive sheet number: 30
